\begin{table}[H]\centering\footnotesize
\caption{$t$-sweep of three methods on the four collapse datasets (AMI).}
\label{tab:collapse-fair}
\setlength{\tabcolsep}{6pt}
\begin{tabular}{@{}ll cccc@{}}
\toprule
Dataset & method & $t{=}1$ & $t{=}2$ & $t{=}3$ & $t{=}4$\\
\midrule
MSRC-v5      & fixed MDT      & $.704$ & $.736$ & $.728$ & $.720$\\
             & imit.\ (Gram)  & $.712$ & $.724$ & $.701$ & $.706$\\
             & end-to-end     & $.533$ & $.309$ & $.277$ & $.239$\\
\midrule
Caltech101-7 & fixed MDT      & $.547$ & $.588$ & $.619$ & $.620$\\
             & imit.\ (Gram)  & $.643$ & $.545$ & $.647$ & $.644$\\
             & end-to-end     & $.234$ & $.094$ & $.070$ & $.085$\\
\midrule
OutdoorScene & fixed MDT      & $.629$ & $.633$ & $.606$ & $.598$\\
             & imit.\ (Gram)  & $.579$ & $.561$ & $.522$ & $.525$\\
             & end-to-end     & $.340$ & $.001$ & $.098$ & $.177$\\
\midrule
UCI          & fixed MDT      & $.846$ & $.834$ & $.830$ & $.840$\\
             & imit.\ (Gram)  & $.841$ & $.810$ & $.840$ & $.828$\\
             & end-to-end     & $.867$ & $.837$ & $.763$ & $.730$\\
\bottomrule
\end{tabular}
\end{table}

\noindent\textbf{Setup.} For diffusion length $t{=}1..4$, three ways to get the embedding.
fixed MDT = best-of-8 length-$t$ operator, plain SVD (no learning); imit.\ (Gram) = a
$\sqrt{\bp}$-Gram encoder trained to reproduce that fixed operator; end-to-end = learn the
operator under the contrastive loss (Scenario 4, 5-seed mean). Four datasets.

\noindent\textbf{Result.} Fixed and imit are stable across $t$ and coincide (imit reproduces the fixed
operator at every length), while end-to-end collapses as $t$ grows (well-connected kNN graphs over-mix
past one hop), except UCI (one dominant view).

\noindent\textbf{Why imit beats end-to-end.} Two reasons. (i)~Fixed vs.\ moving target: imit
matches the Gram of a fixed, known-good operator (the best-of-8 MDT), so its target never drifts;
end-to-end instead learns the operator, so its target $\mathbf G(\theta)$ moves and slides into a
degenerate, spectrum-flattened optimum. (ii)~Conditioning: Gram matching is a well-behaved
quadratic, $\lVert FF^\top-\mathbf G\rVert^2$ with gradient $4(FF^\top-\mathbf G)F$, so it descends
cleanly, whereas the contrastive loss saturates once the operator's rows are near-uniform (its
$\exp$/softmax gives tiny, near-equal gradients) and stalls. imit therefore reaches higher AMI
because it inherits the good fixed operator intact, while end-to-end degrades that operator as it
trains.
